% Part: first-order-logic
% Chapter: axiomatic-deduction
% Section: proving-things

\documentclass[../../../include/open-logic-section]{subfiles}

\begin{document}

\iftag{FOL}
      {\olfileid[pt]{fol}{axd}{pro}}
      {\olfileid[pt]{pl}{axd}{pro}}

\olsection{Exemplos de \usetoken{P}{derivation}}


\begin{ex}
Suponha que queiramos provar $(\lnot !D \lor !E) \lif (!D \lif
!E)$. É claro que essa não é uma instância de nenhum de nossos
axiomas, de modo que precisamos usar a regra \MP{} para !!{derive} essa fórmula.
Nossa única regra é MP, que, dados $!A$ e $!A \lif !B$, permite
justificar~$!B$. Uma estratégia seria usar \olref[prp]{ax:lor3} com
$!A$ igual a $\lnot !D$, $!B$ igual a $!E$ e $!C$ igual a $!D \lif
!E$, isto é, a instância
\[
(\lnot !D \lif (!D \lif !E)) \lif ((!E \lif (!D \lif !E)) \lif ((\lnot
!D \lor !E) \lif (!D \lif !E))).
\]
Por quê? Duas aplicações de MP produzem a última parte, que é o que
queremos. E vemos facilmente que $\lnot !D \lif (!D \lif !E)$ é uma
instância de \olref[prp]{ax:lnot2}, e $!E \lif (!D \lif !E)$ é uma
instância de \olref[prp]{ax:lif1}. Portanto, nossa !!{derivation} é:
\begin{derivation}
  1. &  $\lnot !D \lif (!D \lif !E)$ & \olref[prp]{ax:lnot2} \\
  2. & $(\lnot !D \lif (!D \lif !E)) \lif {}$\\
  &\qquad $((!E \lif (!D \lif !E)) \lif ((\lnot !D \lor !E) \lif (!D \lif !E)))$ & \olref[prp]{ax:lor3}\\
  3. & $(!E \lif (!D \lif !E)) \lif ((\lnot !D \lor !E) \lif (!D \lif !E))$ & 1, 2, \MP\\
  4. & $!E \lif (!D \lif !E)$ & \olref[prp]{ax:lif1}\\
  5. & $(\lnot !D \lor !E) \lif (!D \lif !E)$ & 3, 4, \MP
\end{derivation}
\end{ex}

\begin{ex}\ollabel{ex:identity}
Vamos tentar encontrar !!a{derivation} de $!D \lif !D$. Ela não é uma
instância de um axioma, de modo que precisamos usar \MP{} para
!!{derive} essa fórmula. \olref[prp]{ax:lif1} é um axioma da forma~$!A \lif !B$
ao qual poderíamos aplicar~\MP. Para que fosse útil, é claro, o $!B$
que \MP{} justificaria como um passo correto nesse caso teria de ser
$!D \lif !D$, pois é isso que queremos !!{derive}. Isso significa que
$!A$ também teria de ser $!D$, isto é, poderíamos examinar esta
instância de \olref[prp]{ax:lif1}:
\[
!D \lif (!D \lif !D)
\]
Para aplicar \MP, também precisaríamos justificar a segunda premissa
correspondente, isto é,~$!A$. Mas, em nosso caso, ela seria~$!D$, e não
conseguiremos !!{derive}~$!D$ isoladamente. Portanto, precisamos de uma
estratégia diferente.

O outro axioma que envolve apenas~$\lif$ é \olref[prp]{ax:lif2}, isto é,
\[
(!A \lif (!B \lif !C)) \lif ((!A \lif !B) \lif (!A \lif !C))
\]
Poderíamos chegar à última condicional encaixada aplicando \MP{} duas
vezes. Novamente, isso significaria que queremos uma instância de
\olref[prp]{ax:lif2} na qual $!A \lif !C$ seja $!D \lif !D$, a
!!{formula} que buscamos. Então, é claro, $!A$ e $!C$ são ambos~$!D$.
Como devemos escolher~$!B$ para que tanto $!A \lif (!B \lif !C)$ quanto
$!A \lif !B$, isto é, em nosso caso $!D \lif (!B \lif !D)$ e $!D \lif
!B$, sejam também !!{derivable}? Bem, a primeira já é uma instância de
\olref[prp]{ax:lif1}, seja qual for nossa escolha de $!B$. E $!D \lif
!B$ seria outra instância de \olref[prp]{ax:lif1} se $!B$ fosse $(!D
\lif !D)$. Portanto, nossa !!{derivation} é:
\begin{derivation}
  1. & $!D \lif ((!D \lif !D) \lif !D)$ & \olref[prp]{ax:lif1}\\
  2. & $(!D \lif ((!D \lif !D) \lif !D)) \lif {}$\\
  & \qquad $((!D \lif (!D \lif !D)) \lif (!D \lif !D))$ & \olref[prp]{ax:lif2}\\
  3. & $(!D \lif (!D \lif !D)) \lif (!D \lif !D)$ & 1, 2, \MP\\
  4. & $!D \lif (!D \lif !D)$ & \olref[prp]{ax:lif1}\\
  5. & $!D \lif !D$ & 3, 4, \MP
\end{derivation}
\end{ex}

\begin{ex}\ollabel{ex:chain}
Às vezes, queremos mostrar que existe !!a{derivation} de alguma
!!{formula} a partir de outras !!{formula}s~$\Gamma$. Por exemplo,
mostremos que podemos !!{derive} $!A \lif !C$ a partir de $\Gamma =
\{!A \lif !B, !B \lif !C\}$.
\begin{derivation}
  1. & $!A \lif !B$ & \Hyp\\
  2. & $!B \lif !C$ & \Hyp\\
  3. & $(!B \lif !C) \lif (!A \lif (!B \lif !C))$ & \olref[prp]{ax:lif1} \\
  4. & $!A \lif (!B \lif !C)$ & 2, 3, \MP\\
  5. & $(!A \lif (!B \lif !C)) \lif {}$\\
  & \qquad $((!A \lif !B) \lif (!A \lif !C))$ & \olref[prp]{ax:lif2}\\
  6. & $((!A \lif !B) \lif (!A \lif !C))$ & 4, 5, \MP\\
  7. & $!A \lif !C$ & 1, 6, \MP
\end{derivation}
As linhas rotuladas ``\Hyp'' (de ``hipótese'') indicam que a !!{formula}
nessa linha é !!a{element} de~$\Gamma$.
\end{ex}

\begin{prop}
  \ollabel{prop:chain} Se $\Gamma \Proves !A \lif !B$ e $\Gamma
  \Proves !B \lif !C$, então $\Gamma \Proves !A \lif !C$
\end{prop}

\begin{proof}
  Suponha que $\Gamma \Proves !A \lif !B$ e $\Gamma \Proves !B \lif
  !C$. Então existe !!a{derivation} de $!A \lif !B$ a partir
  de~$\Gamma$; e também existe !!a{derivation} de~$!B \lif !C$ a partir
  de~$\Gamma$. Combine-as em uma única !!{derivation}, concatenando-as.
  Agora acrescente as linhas 3--7 da !!{derivation} do exemplo anterior.
  Isso é !!a{derivation} de $!A \lif !C$---que é a última linha da nova
  !!{derivation}---a partir de~$\Gamma$. Observe que as justificativas
  das linhas 4 e~7 continuam válidas se a referência à linha número~1
  for substituída por uma referência à última linha da !!{derivation}
  de~$!A \lif !B$, e a referência à linha número~2 por uma referência à
  última linha da !!{derivation} de~$!B \lif !C$.
\end{proof}

\begin{prob} 
Mostre que as afirmações a seguir valem exibindo !!{derivation}s a
partir dos axiomas:
\begin{enumerate} 
\item $(!A \land !B) \lif (!B \land !A)$
\item $((!A \land !B) \lif !C) \lif (!A \lif (!B \lif !C))$
\item $\lnot(!A \lor !B) \lif \lnot !A$
\end{enumerate} 
\end{prob}



\end{document}
